% GENERATED FILE. Edit canonical lessons, then run npm run generate:books.
% canonical-source-hash: fnv1a64:460dfe05608e8d3c
% canonical-lessons: ES-C772-brisa, ES-C772-soplar, ES-C772-ladera, ES-C772-resbalar, ES-C772-brujula

\chapter{Cuesta arriba}
\label{ch:es-cuesta-arriba}
\hlchaptermodality{772}

\begin{chapteropening}
\textbf{By the end of this chapter:} \emph{I can cope on a hill walk: a breeze blowing, a steep slope, slipping, and finding the way with a compass.}

Complete the last lesson of chapter 772: I can cope on a hill walk: a breeze blowing, a steep slope, slipping, and finding the way with a compass.
\end{chapteropening}

\section[la brisa]{la brisa — the breeze}

\label{lesson:ES-C772-brisa}



\begin{quote}
Before the new one: say the Spanish for raw, then the Spanish for to smell.
\end{quote}

\subsection*{You'll want to know: la brisa}
\textbf{la brisa} — \textquotedblleft{}the breeze\textquotedblright{}.

\begin{quote}
Arriba sopla una brisa fresca.
\end{quote}

\begin{quote}
\emph{Up there a cool breeze blows.}
\end{quote}

\emph{La brisa del mar} — the sea breeze.

\begin{cousinweb}
\emph{brisa} comes from old Spanish \emph{briza} (EtymDB), which named the northeast wind. Where \emph{briza} came from is not settled; French \emph{brise} and English \emph{breeze} are related words.
\end{cousinweb}

\begin{grammarlens}[title={What you've built}]
One more word for telling, explaining and describing at B1.
\end{grammarlens}

\subsection*{Your turn}
\begin{itemize}
\raggedright
  \item \emph{Say it:} \emph{la brisa}
  \item \emph{Say it:} \emph{la brisa}, once more
  \item \emph{Say it:} a sentence of your own with \emph{la brisa}
  \item \emph{Recall:} say the Spanish for raw, then the Spanish for to smell, then say \emph{la brisa} again
\end{itemize}

\begin{tcolorbox}[breakable,colback=teal!4,colframe=teal!35!black,title={Before you move on}]
What does la brisa mean? (\textquotedblleft{}The breeze\textquotedblright{}.) Say it once more.
\end{tcolorbox}
\section[soplar]{soplar — to blow}

\label{lesson:ES-C772-soplar}



\begin{quote}
Before the new one: say the Spanish for to smell, then the Spanish for the breeze.
\end{quote}

\subsection*{You'll want to know: soplar}
\textbf{soplar} — \textquotedblleft{}to blow\textquotedblright{}.

\begin{quote}
Si sopla mucho viento, volvemos.
\end{quote}

\begin{quote}
\emph{If the wind blows hard, we'll turn back.}
\end{quote}

\emph{Soplar las velas} — to blow out the candles.

\begin{cousinweb}
\emph{soplar} is inherited from Latin \emph{sufflāre} 'to blow from below, to puff out' (Lewis \& Short), built on \emph{sub} and \emph{flāre} 'to blow'.
\end{cousinweb}

\begin{grammarlens}[title={What you've built}]
One more word for telling, explaining and describing at B1.
\end{grammarlens}

\subsection*{Your turn}
\begin{itemize}
\raggedright
  \item \emph{Say it:} \emph{soplar}
  \item \emph{Say it:} \emph{soplar}, once more
  \item \emph{Say it:} a sentence of your own with \emph{soplar}
  \item \emph{Recall:} say the Spanish for to smell, then the Spanish for the breeze, then say \emph{soplar} again
\end{itemize}

\begin{tcolorbox}[breakable,colback=teal!4,colframe=teal!35!black,title={Before you move on}]
What does soplar mean? (\textquotedblleft{}To blow\textquotedblright{}.) Say it once more.
\end{tcolorbox}
\section[la ladera]{la ladera — the hillside, the slope}

\label{lesson:ES-C772-ladera}



\begin{quote}
Before the new one: say the Spanish for the breeze, then the Spanish for to blow.
\end{quote}

\subsection*{You'll want to know: la ladera}
\textbf{la ladera} — \textquotedblleft{}the hillside, the slope\textquotedblright{}.

\begin{quote}
El camino sube por la ladera de la montaña.
\end{quote}

\begin{quote}
\emph{The path climbs up the side of the mountain.}
\end{quote}

\begin{cousinweb}
\emph{ladera} is built on \emph{lado} 'side' with the ending \emph{-era}. \emph{Lado} is inherited from Latin \emph{latus} 'side, flank' (EtymDB).
\end{cousinweb}

\begin{grammarlens}[title={What you've built}]
One more word for telling, explaining and describing at B1.
\end{grammarlens}

\subsection*{Your turn}
\begin{itemize}
\raggedright
  \item \emph{Say it:} \emph{la ladera}
  \item \emph{Say it:} \emph{la ladera}, once more
  \item \emph{Say it:} a sentence of your own with \emph{la ladera}
  \item \emph{Recall:} say the Spanish for the breeze, then the Spanish for to blow, then say \emph{la ladera} again
\end{itemize}

\begin{tcolorbox}[breakable,colback=teal!4,colframe=teal!35!black,title={Before you move on}]
What does la ladera mean? (\textquotedblleft{}The hillside, the slope\textquotedblright{}.) Say it once more.
\end{tcolorbox}
\section[resbalar]{resbalar — to slip}

\label{lesson:ES-C772-resbalar}



\begin{quote}
Before the new one: say the Spanish for to blow, then the Spanish for the hillside.
\end{quote}

\subsection*{You'll want to know: resbalar}
\textbf{resbalar} — \textquotedblleft{}to slip\textquotedblright{}.

\begin{quote}
Cuidado, que las piedras resbalan.
\end{quote}

\begin{quote}
\emph{Careful, the stones are slippery.}
\end{quote}

\emph{Resbalarse} is to slip and fall: \emph{Me resbalé en el barro.}

\begin{cousinweb}
\emph{resbalar} is built on old Spanish \emph{esbarar} 'to slide' with \emph{re-} (EtymDB). Where \emph{esbarar} came from is not settled.
\end{cousinweb}

\begin{grammarlens}[title={What you've built}]
One more word for telling, explaining and describing at B1.
\end{grammarlens}

\subsection*{Your turn}
\begin{itemize}
\raggedright
  \item \emph{Say it:} \emph{resbalar}
  \item \emph{Say it:} \emph{resbalar}, once more
  \item \emph{Say it:} a sentence of your own with \emph{resbalar}
  \item \emph{Recall:} say the Spanish for to blow, then the Spanish for the hillside, then say \emph{resbalar} again
\end{itemize}

\begin{tcolorbox}[breakable,colback=teal!4,colframe=teal!35!black,title={Before you move on}]
What does resbalar mean? (\textquotedblleft{}To slip\textquotedblright{}.) Say it once more.
\end{tcolorbox}
\section[la brújula]{la brújula — the compass}

\label{lesson:ES-C772-brujula}



\begin{quote}
Before the new one: say the Spanish for the hillside, then the Spanish for to slip.
\end{quote}

\subsection*{You'll want to know: la brújula}
\textbf{la brújula} — \textquotedblleft{}the compass\textquotedblright{}.

\begin{quote}
Sin mapa, usamos la brújula.
\end{quote}

\begin{quote}
\emph{With no map, we used the compass.}
\end{quote}

\begin{cousinweb}
\emph{brújula} comes from Italian \emph{bussola} 'a compass' (EtymDB), from Late Latin \emph{buxula} 'a little box', built on \emph{buxus} 'boxwood' (Lewis \& Short). The needle sat in a small wooden box.
\end{cousinweb}

\begin{grammarlens}[title={What you've built}]
One more word for telling, explaining and describing at B1.
\end{grammarlens}

\subsection*{Your turn}
\begin{itemize}
\raggedright
  \item \emph{Say it:} \emph{la brújula}
  \item \emph{Say it:} \emph{la brújula}, once more
  \item \emph{Say it:} a sentence of your own with \emph{la brújula}
  \item \emph{Recall:} say the Spanish for the hillside, then the Spanish for to slip, then say \emph{la brújula} again
\end{itemize}

\begin{tcolorbox}[breakable,colback=teal!4,colframe=teal!35!black,title={Before you move on}]
What does la brújula mean? (\textquotedblleft{}The compass\textquotedblright{}.) Say it once more.
\end{tcolorbox}
